% ============ §22 The local flatness criterion (原书 pp.173-178) ============
\section{The local flatness criterion}\label{sec:22}\index{local flatness criterion, section 22}\index{flatness, local criterion, section 22}

\begin{theorem}\label{thm:22.1}
Let $A$ be a ring, $B$ a Noetherian $A$-algebra, $M$ a finite $B$-module, and
$J$ an ideal of $B$ contained in $\rad(B)$; set $M_n = M/J^{n+1}M$ for $n \ge 0$.
If $M_n$ is flat over $A$ for every $n \ge 0$, then $M$ is also flat over $A$.
\end{theorem}

\begin{proof}
According to Theorem~\ref{thm:7.7}, we need only show that for a finitely
generated ideal $I$ of $A$, the standard map $u\colon I \otimes_A M \to M$ is
injective. Set $I \otimes_A M = M'$; then $M'$ is also a finite $B$-module, and
hence is separated for the $J$-adic topology. Let $x \in \Ker(u)$; we prove that
$x \in \cap_n J^nM' = 0$. For each $n$ we have $M'/J^{n+1}M' = (I \otimes_A M)
\otimes_B B/J^{n+1} = I \otimes_A M_n$, and the standard map $I \otimes_A M_n \to
M_n$ is injective, by the assumption that $M_n$ is flat. Then we deduce that
$x \in J^{n+1}M'$ from the commutative diagram
\[
\begin{tikzcd}
M' = I \otimes_A M \arrow[r] \arrow[d, "\bar{u}"'] & I \otimes_A M_n \arrow[d, "u_n"] \\
M \arrow[r] & M_n
\end{tikzcd}
\]
in which the horizontal maps are the natural ones. Indeed, the image of $x$ in
$I \otimes_A M_n$ is mapped to $0$ in $M_n$, hence is $0$, so that
$x \in J^{n+1}M'$. Since this holds for all $n$, $x \in \cap_n J^nM' = 0$ (the
intersection being zero by the Krull intersection theorem, because $J \subset
\rad(B)$ and $M'$ is a finite $B$-module). Thus $u$ is injective, and $M$ is
flat over $A$.
\end{proof}

\begin{theorem}\label{thm:22.2}
Let $A$ be a ring, $B$ a Noetherian $A$-algebra, and $M$ a finite $B$-module;
suppose that $b$ is an $M$-regular element of $\rad(B)$. Then if $M/bM$ is flat
over $A$, so is $M$.
\end{theorem}

\begin{proof}
For each $i \ge 0$ the sequence $0 \to M/b^iM \xrightarrow{b} M/b^{i+1}M \to
M/bM \to 0$ is exact, so that by Theorem~\ref{thm:7.9} and an induction on $i$,
every $M/b^iM$ is flat over $A$. Thus we can just apply the previous theorem.
\end{proof}

\begin{definition*}
Let $A$ be a ring and $I$ an ideal of $A$; an $A$-module $M$ is said to be
\emph{$I$-adically ideal-separated}\index{I-adically ideal-separated, section 22} if $a \otimes_A M$ is separated for the $I$-adic
topology for every finitely generated ideal $a$ of $A$.

For example, if $B$ is a Noetherian $A$-algebra and $IB \subset \rad(B)$ then a
finite $B$-module $M$ is $I$-adically ideal-separated as an $A$-module.
\end{definition*}

Let $A$ be a ring, $I$ an ideal of $A$ and $M$ an $A$-module. Set $A_n = A/I^{n+1}$,
$M_n = M/I^{n+1}M$ for $n \ge 0$ and $\gr(A) = \bigoplus_{n \ge 0} I^n/I^{n+1}$,
$\gr(M) = \bigoplus_{n \ge 0} I^nM/I^{n+1}M$. There exist standard maps
\[
  \gamma_n\colon (I^n/I^{n+1}) \otimes_{A_0} M_0 \to I^nM/I^{n+1}M \quad \text{for } n \ge 0,
\]
and we can put together the $\gamma_n$ into a morphism of $\gr(A)$-modules
\[
  \gamma\colon \gr(A) \otimes_{A_0} M_0 \to \gr(M).
\]

\begin{theorem}\label{thm:22.3}
In the above notation, suppose that one of the following two conditions is
satisfied:
\begin{enumerate}[(a)]
  \item $I$ is a nilpotent ideal;
or ($\beta$) $A$ is a Noetherian ring and $M$ is $I$-adically ideal-separated.
\end{enumerate}
Then the following conditions are equivalent.
\begin{enumerate}[(1)]
  \item $M$ is flat over $A$;
  \item $\Tor^A_1(N, M) = 0$ for every $A_0$-module $N$;
  \item $M_0$ is flat over $A_0$, and $I \otimes_A M = IM$;
  \item[$3'$] $M_0$ is flat over $A_0$, and $\Tor^A_1(A_0, M) = 0$;
  \item $M_0$ is flat over $A_0$, and $\gamma_n$ is an isomorphism for every $n \ge 0$;
  \item[$4'$] $M_0$ is flat over $A_0$, and $\gamma$ is an isomorphism;
  \item $M_n$ is flat over $A_n$ for every $n \ge 0$.
\end{enumerate}
In fact, the implications (1) $\Rightarrow$ (2) $\Leftarrow$ (3) $\Leftarrow$ (3$'$) $\Leftarrow$ (4) $\Leftarrow$ (5)
hold without any assumption on $M$.
\end{theorem}

\begin{proof}
First of all, let $M$ be arbitrary.

(1) $\Rightarrow$ (2) is trivial.

(2) $\Rightarrow$ (3) If $N$ is an $A_0$-module then we have
\[
  N \otimes_A M = (N \otimes_{A_0} A_0) \otimes_A M = N \otimes_{A_0} M_0,
\]
and hence for an exact sequence $0 \to N_1 \to N_2 \to N_3 \to 0$ of
$A_0$-modules we get an exact sequence
\[
  0 = \Tor^A_1(N_3, M) \to N_1 \otimes_{A_0} M_0 \to N_2 \otimes_{A_0} M_0 \to
      N_3 \otimes_{A_0} M_0 \to 0;
\]
therefore $M_0$ is flat over $A_0$. Also, from the exact sequence $0 \to I \to
A \to A_0 \to 0$ we get an exact sequence
\[
  0 = \Tor^A_1(A_0, M) \to I \otimes_A M \to M \to M_0 \to 0,
\]
so that $I \otimes_A M = IM$.

(3) $\Rightarrow$ (3$'$) is clear, and (3$'$) $\Rightarrow$ (2) follows as follows: if
$N$ is an $A_0$-module, we can choose an exact sequence of $A_0$-modules
$0 \to R \to F_0 \to N \to 0$ with $F_0$ a free $A_0$-module. From this we get
the exact sequence
\[
  \Tor^A_1(F_0, M) = 0 \to \Tor^A_1(N, M) \to R \otimes_{A_0} M_0 \to F_0 \otimes_{A_0} M_0,
\]
and since $M_0$ is flat over $A_0$, the final arrow is injective, so that
$\Tor^A_1(N, M) = 0$.

(3) $\Rightarrow$ (4) We have already seen that $\gamma_0$ is the identity. For $n = 1$,
from the exact sequence $0 \to I^2 \to I \to I/I^2 \to 0$ we get
$\Tor^A_1(I/I^2, M) = 0$, so that from $0 \to I^2 \to I \to I/I^2 \to 0$ the sequence
$0 \to I^2 \otimes_A M \to I \otimes_A M \to (I/I^2) \otimes_A M \to 0$
is exact. From $I \otimes_A M = IM$ we get $I^2 \otimes_A M = I^2M$ and
$(I/I^2) \otimes_A M \cong IM/I^2M$. Proceeding similarly, from
$0 \to I^{n+1} \to I^n \to I^n/I^{n+1} \to 0$ we get by induction
$I^{n+1} \otimes_A M = I^{n+1}M$ and $(I^n/I^{n+1}) \otimes_A M \cong I^nM/I^{n+1}M$.
(4$'$) is just a restatement of (4).

(4) $\Rightarrow$ (5) We fix an $n > 0$ and prove that $M_n$ is flat over $A_n$. For
$i \le n$ we have a commutative diagram
\[
\begin{tikzcd}
0 \arrow[r] &
(I^{i+1}/I^{n+1}) \otimes_A M \arrow[r] \arrow[d, "\alpha_{i+1}"] &
(I^i/I^{n+1}) \otimes_A M \arrow[r] \arrow[d, "\alpha_i"] &
(I^i/I^{i+1}) \otimes_A M \arrow[r] \arrow[d, "\gamma_i"] & 0 \\
0 \arrow[r] &
I^{i+1}M/I^{n+1}M \arrow[r] &
I^iM/I^{n+1}M \arrow[r] &
I^iM/I^{i+1}M \arrow[r] & 0
\end{tikzcd}
\]
with exact rows (the top row is exact since $\Tor^A_1(I^i/I^{i+1}, M) = 0$). By
assumption $\gamma_i$ is an isomorphism, and since $\alpha_{i+1}$ is an isomorphism,
by downwards induction on $i$ we see that $\alpha_n, \ldots, \alpha_1$ are
isomorphisms. In particular,
\[
  \alpha_1\colon (I/I^{n+1}) \otimes_A M \xrightarrow{\cong} IM/I^{n+1}M,
\]
and $(I/I^{n+1}) \otimes_A M = I \otimes_{A_n} M_n = IM_n$,
so that the conditions in (3) are satisfied by $A_n$, $M_n$ and $I/I^{n+1}$.
Therefore by (3) $\Rightarrow$ (2), we have $\Tor^{A_n}_1(N, M_n) = 0$ for every
$A_0$-module $N$. Now if $N$ is an $A_i$-module then $IN$ and $N/IN$ are both
$A_{i-1}$-modules, and
\[
  0 \to IN \to N \to N/IN \to 0
\]
is exact, so that by induction on $i$ we get finally that $\Tor^{A_n}_1(N, M_n) = 0$
for all $A_n$-modules $N$. Therefore $M_n$ is a flat $A_n$-module.

Next, assuming either (a) or ($\beta$) we prove (5) $\Rightarrow$ (1). In case (a) we have
$A = A_n$ and $M = M_n$ for large enough $n$, so that this is clear. In case
($\beta$), let $a$ be a finitely generated ideal of $A$; we prove that the standard map
$j\colon a \otimes_A M \to M$ is injective. By hypothesis we have
$\cap_n I^n(a \otimes_A M) = 0$, so that we need only prove that $\Ker(j) \subset
I^n(a \otimes_A M)$ for all $n > 0$. For a fixed $n$, by the Artin-Rees lemma,
$I^k \cap a \subset I^na$ for sufficiently large $k > n$. We now consider the
natural map
\[
  a \otimes_A M \to (a/I^ka) \otimes_A M = (a \otimes_A M)/I^k(a \otimes_A M).
\]
Since $M_{k-1}$ is flat over $A_{k-1} = A/I^k$, the map
\[
  (a/I^ka) \otimes_A M = (a/I^ka) \otimes_{A_{k-1}} M_{k-1} \to M_{k-1}
\]
is injective, so that from the commutative diagram
\[
\begin{tikzcd}
a \otimes_A M \arrow[r] \arrow[d, "j"'] & (a/I^ka) \otimes_A M \arrow[d] \\
M \arrow[r] & M_{k-1}
\end{tikzcd}
\]
we get $\Ker(j) \subset \Ker(f) \subset \Ker(gf) = I^n(a \otimes_A M)$. This is what
we needed to prove.
\end{proof}

This theorem is particularly effective when $A$ is a Noetherian local ring and
$I$ is the maximal ideal, since if $A_0$ is a field, $M_0$ is automatically flat
over $A_0$ in (3)---(4$'$). Also, in this case, requiring $M_n$ to be flat over $A_n$ in
(5) is the same as requiring it to be a free $A_n$-module, by
Theorem~\ref{thm:7.10}.

We now discuss some applications of the above theorem.

\begin{theorem}\label{thm:22.4}
Let $(A, m)$ and $(B, n)$ be Noetherian local rings, $\hat{A}$ and $\hat{B}$
their respective completions, and $A \to B$ a local homomorphism.
\begin{enumerate}[(i)]
  \item For $M$ a finite $B$-module, set $\hat{M} = M \otimes_B \hat{B}$; then
        \[
          M \text{ is flat over } A \iff \hat{M} \text{ is flat over } \hat{A} \iff \hat{M} \text{ is flat over } A.
        \]
  \item Writing $M^*$ for the $(mB)$-adic completion of $M$ we have
        \[
          M \text{ is flat over } A \iff M^* \text{ is flat over } \hat{A} \iff M^* \text{ is flat over } A.
        \]
\end{enumerate}
\end{theorem}

\begin{proof}
(i) The first equivalence comes from the transitivity law for flatness, together
with the fact that $\hat{B}$ is faithfully flat over $B$; the second, from the fact that
both sides are equivalent to $\hat{B}/m^n\hat{B}$ being flat over $A/m^n$ for all $n > 0$.

(ii) All three conditions are equivalent to $M/m^nM$ being flat over $A/m^n$ for
all $n$.
\end{proof}

\begin{theorem}\label{thm:22.5}
Let $(A, m, k)$ and $(B, n, k')$ be Noetherian local rings, $A \to B$ a local
homomorphism, and $u\colon M \to N$ a morphism of finite $B$-modules. Then if
$N$ is flat over $A$, the following two conditions are equivalent:
\begin{enumerate}[(1)]
  \item $u$ is injective and $N/u(M)$ is flat over $A$;
  \item $\bar{u}\colon M \otimes_A k \to N \otimes_A k$ is injective.
\end{enumerate}
\end{theorem}

\begin{proof}
(1) $\Rightarrow$ (2) is easy, so we only give the proof of (2) $\Rightarrow$ (1). Suppose that
$x \in M$ is such that $u(x) = 0$; then $\bar{u}(\bar{x}) = 0$, so that $\bar{x} = 0$, in
other words, $x \in mM$. Now assuming $x \in m^nM$, we will deduce $x \in m^{n+1}M$.
Let $\{a_1, \ldots, a_r\}$ be a minimal basis of the $A$-module $m^n$, and write
$x = \sum a_i y_i$ with $y_i \in M$; then $0 = \sum a_i u(y_i)$. Since $N$ is flat over
$A$, by Theorem~\ref{thm:7.6} there exist $c_{ij} \in A$ and $z_j \in N$ such that
\[
  \sum_i a_i c_{ij} = 0 \quad \text{for all } j, \quad \text{and} \quad u(y_i) = \sum_j c_{ij}z_j \quad \text{for all } i.
\]
By choice of $a_1, \ldots, a_r$, all the $c_{ij} \in m$, and hence $u(y_i) \in mN$ and
$\bar{u}(\bar{y}_i) = 0$, so that $\bar{y}_i = 0$, and $y_i \in mM$. Therefore
$x \in m^{n+1}M$. We have proved that $x \in \cap_n m^nM = 0$, and hence $u$ is
injective. Now from $0 \to M \xrightarrow{u} N \to N/u(M) \to 0$ we get
$\Tor^A_1(k, N/u(M)) = 0$, so that by Theorem~\hyperref[thm:22.3]{3},
$N/u(M)$ is flat over $A$.
\end{proof}

\begin{corollary*}
Let $A$, $B$ and $A \to B$ be as above, and $M$ a finite $B$-module; set
$\bar{B} = B \otimes_A k = B/mB$, and for $x_1, \ldots, x_n \in n$ write $\bar{x}_i$ for the
images in $\bar{B}$ of $x_i$. Then the following conditions are equivalent:
\begin{enumerate}[(1)]
  \item $x_1, \ldots, x_n$ is an $M$-sequence and $M_n = M/\sum x_iM$ is flat over $A$;
  \item $\bar{x}_1, \ldots, \bar{x}_n$ is an $M \otimes_A k$-sequence and $M$ is flat over $A$.
\end{enumerate}
\end{corollary*}

\begin{proof}
(2) $\Rightarrow$ (1) follows at once from the theorem. For (1) $\Rightarrow$ (2) we must prove
that $M_i = M/(x_1M + \cdots + x_iM)$ is flat for $i = 1, \ldots, n$; but if $M_i$ is flat
over $A$ then by Theorem~\hyperref[thm:22.2]{2}, so is $M_{i-1}$.
\end{proof}

\begin{theorem}\label{thm:22.6}
Let $A$ be a Noetherian ring, $B$ a Noetherian $A$-algebra, $M$ a finite
$B$-module, and $b \in B$ a given element. Suppose that $M$ is flat over $A$ and
that $b$ is $M/(P \cap A)$-regular for every maximal ideal $P$ of $B$; then $b$ is
$M$-regular and $M/bM$ is flat over $A$.
\end{theorem}

\begin{proof}
Write $K$ for the kernel of $M \xrightarrow{b} M$; then $K = \cap_P K_P = 0$ for all $P$.
Hence $b$ is $M$-regular if and only if $b$ is $M_P$-regular for all $P$. Moreover,
according to Theorem~\ref{thm:7.1}, $A$-flatness is also a local property in both $A$
and $B$, so that we can replace $B$ by $B_P$ (for a maximal ideal $P$ of $B$), $A$ by
$A_{(P \cap A)}$ and $M$ by $M_P$, and this case reduces to Theorem~\hyperref[thm:22.5]{5}.
\end{proof}

\begin{corollary*}
Let $A$ be a Noetherian ring, $B = A[X_1, \ldots, X_n]$ the polynomial ring over
$A$, and let $f \in B$. If the ideal of $A$ generated by the coefficients of $f$
contains $1$ then $f$ is a non-zero-divisor of $B$, and $B/fB$ is flat over $A$.
The same thing holds for the formal power series ring $B = A[[X_1, \ldots, X_n]]$.
\end{corollary*}

\begin{proof}
The polynomial ring is a free $A$-module, and therefore flat; the formal power
series ring is flat by Ex.~\ref{ex:7.4-ex}. Furthermore, for $\mathfrak{p} \in \Spec A$, if
$B = A[X_1, \ldots, X_n]$ then $B/\mathfrak{p}B = (A/\mathfrak{p})[X_1, \ldots, X_n]$, and in the formal
power series case we also have $B/\mathfrak{p}B = (A/\mathfrak{p})[[X_1, \ldots, X_n]]$ since
$\mathfrak{p}$ is finitely generated. In either case $B/\mathfrak{p}B$ is an integral domain,
so that the assertion follows directly from the theorem.
\end{proof}

\begin{remark}[\textbf{Flatness of a graded module}]\index{graded module, flatness of, section 22}
Let $G$ be an Abelian group, $R = \bigoplus_{g \in G} R_g$ a $G$-graded ring and
$M = \bigoplus_{g \in G} M_g$ a graded $R$-module, not necessarily finitely
generated.
\begin{enumerate}[(1)]
  \item The following three conditions are equivalent:
        \begin{enumerate}[(a)]
          \item $M$ is $R$-flat;
          \item If $\mathcal{F}\colon \cdots \to N \to N' \to N'' \to \cdots$ is an exact
                sequence of graded $R$-modules and $R$-linear maps preserving
                degrees, then $\mathcal{F} \otimes_R M$ is exact;
          \item $\Tor^R_1(M, R/H) = 0$ for every finitely generated homogeneous
                ideal $H$ of $R$.
        \end{enumerate}
        The proof is left to the reader as an exercise, or can be found in
        Herrmann and Orbanz~[3]. Using this criterion one can adapt the proof of
        Theorem~\hyperref[thm:22.3]{3} to prove the following graded version.
  \item Let $I$ be a (not necessarily homogeneous) ideal of $R$. Suppose that
        \begin{enumerate}[(i)]
          \item for every finitely generated homogeneous ideal $H$ of $R$, the
                $R$-module $H \otimes_R M$ is $I$-adically separated;
          \item $M_0 = M/IM$ is $R/I$-flat;
          \item $\Tor^R_1(M, R/I) = 0$.
        \end{enumerate}
        Then $M$ is $R$-flat.
  \item As an application one can prove the following:
\end{enumerate}
\begin{enumerate}[(3)]
  \item Let $A = \bigoplus_{n \ge 0} A_n$ and $B = \bigoplus_{n \ge 0} B_n$ be graded Noetherian
        rings. Assume that $A_0$, $B_0$ are local rings with maximal ideals $m$, $n$ and
        set $M = m + A_1 + A_2 + \cdots$, $N = n + B_1 + B_2 + \cdots$; let
        $f\colon A \to B$ be a ring homomorphism of degree $0$ such that $f(m) \subset n$.
        Then the following are equivalent:
        \begin{enumerate}[(a)]
          \item $B$ is $A$-flat;
          \item $B_N$ is $A_M$-flat;
          \item $B_0$ is $A_0$-flat.
        \end{enumerate}
\end{enumerate}
\end{remark}

% ============ 习题 ============
\section*{Exercises to \S\ref{sec:22}}
\emph{Prove the following propositions.}

\begin{exercise}\label{ex:22.1-ex}
(The Nagata flatness theorem, see [N1], p.~65). Let $(A, m, k)$ and $(B, n, k')$
be Noetherian local rings, and suppose that $A \subset B$ and that $mB$ is an
$n$-primary ideal. We say that the \emph{transition theorem} holds between $A$ and
$B$ if $l(A/q) \cdot l(B/mB) = l(B/qB)$ for every $m$-primary ideal $q$ of $A$. This
holds if and only if $B$ is flat over $A$.
\end{exercise}

\begin{exercise}\label{ex:22.2-ex}
Let $(A, m)$ be a Noetherian local ring, and $k \subset A$ a subfield. If
$x_1, \ldots, x_n \in m$ is an $A$-sequence then $x_1, \ldots, x_n$ are algebraically
independent over $k$, and $A$ is flat over $C = k[x_1, \ldots, x_n]$
(Hartshorne [2]).
\end{exercise}

\begin{exercise}\label{ex:22.3-ex}
Let $(A, m, k)$ be a Noetherian local ring, $B$ a Noetherian $A$-algebra, and $M$
a finite $B$-module. Suppose that $mB \subset \rad(B)$. If $x \in m$ is both
$A$-regular and $M$-regular, and if $M/xM$ is flat over $A/xA$ then $M$ is flat
over $A$.
\end{exercise}

\begin{exercise}\label{ex:22.4-ex}
Let $A$ be a Noetherian ring and $B$ a flat Noetherian $A$-algebra; if $I$ and $J$
are ideals of $A$ and $B$ such that $IB \subset J$ then the $J$-adic completion of
$B$ is flat over the $I$-adic completion of $A$.
\end{exercise}
